% SOURCE_AUTHORITY: SGA5_LNM589_SCAN; SHA256 B256EBD072A8C68209518412A263C9289C6F1854A346733D86F885930D5FE6CA.
% SOURCE_PAGES: PDF 8 = printed VIII (bibliography [1]--[5]); PDF 9 = printed IX (Liste des SGA); PDF 10 = unnumbered blank corresponding to Roman X; PDF 11--12 = printed XI--XII (Table des matieres, absence note, and dependency diagram).
\clearpage
\pagenumbering{Roman}
\setcounter{page}{8}

% Printed page VIII: the five bibliography entries closing the introduction.
\begin{enumerate}[label={\textup{[\arabic*]}},leftmargin=3.2em,labelsep=1em,itemsep=1.1em,parsep=0pt,topsep=0pt]
\item Deligne, P., \emph{La conjecture de Weil II}, a ser publicado.

\item Grothendieck, A., \emph{Formule de Lefschetz et rationalit\'e des fonctions $L$}, S\'eminaire Bourbaki 64/65, n\textsuperscript{o}~279, em \emph{Dix expos\'es sur la cohomologie des sch\'emas}, North Holland, Pub. Co., 1968.

\item Raynaud, M., \emph{Caract\'eristique d'Euler--Poincar\'e d'un faisceau et cohomologie des vari\'et\'es ab\'eliennes (d'apr\`es Ogg--Shafarevitch et Grothendieck)}, S\'eminaire Bourbaki 64/65, n\textsuperscript{o}~286, em \emph{Dix expos\'es sur la cohomologie des sch\'emas}, North Holland Pub. Co., 1968.

\item Serre, J.-P., \emph{Repr\'esentations lin\'eaires des groupes finis}, Hermann, 2.\textsuperscript{a} edição revista, 1971.

\item Stallings, J., \emph{Centerless groups -- an algebraic formulation of Gottlieb's theorem}, \emph{Topology}, 4, p.~129--134, 1965.
\end{enumerate}

\clearpage
% Printed page IX.
\section*{\uline{Lista das SGA}}

\noindent SGA = Seminário de Geometria Algébrica do Bois-Marie.

\begingroup
\newcommand{\sgalistentry}[2]{%
  \par\medskip\noindent
  \hangindent=4.7em\hangafter=1
  \textbf{#1:}\ #2\par}
\sgalistentry{SGA 1}{\emph{Recobrimentos étales e grupo fundamental}, por A. Grothendieck, \emph{Lecture Notes in Mathematics} n\textsuperscript{o}~224, Springer-Verlag, 1971.}

\sgalistentry{SGA 2}{\emph{Cohomologia local dos feixes coerentes e teoremas de Lefschetz locais e globais}, por A. Grothendieck, North Holland Pub. Co., 1968.}

\sgalistentry{SGA 3}{\emph{Esquemas em grupos I, II, III}, por M. Demazure e A. Grothendieck, \emph{Lecture Notes in Mathematics} n\textsuperscript{os}~151, 152, 153, Springer-Verlag, 1970.}

\sgalistentry{SGA 4}{\emph{Teoria dos topos e cohomologia étale dos esquemas}, por M. Artin, A. Grothendieck e J.L. Verdier, \emph{Lecture Notes in Mathematics} n\textsuperscript{os}~269, 270, 305, Springer-Verlag, 1972--73.}

\sgalistentry{SGA $4^{1/2}$}{\emph{Cohomologia étale}, por P. Deligne, \emph{Lecture Notes in Mathematics} n\textsuperscript{o}~569, Springer-Verlag, 1977.}

\sgalistentry{SGA 5}{\emph{Cohomologia $\ell$-ádica e funções $L$}, por A. Grothendieck, \emph{Lecture Notes in Mathematics} n\textsuperscript{o}~589, Springer-Verlag, 1977.}

\sgalistentry{SGA 6}{\emph{Teoria das interseções e teorema de Riemann--Roch}, por P. Berthelot, A. Grothendieck e L. Illusie, \emph{Lecture Notes in Mathematics} n\textsuperscript{o}~225, Springer-Verlag, 1971.}

\sgalistentry{SGA 7}{\emph{Grupos de monodromia em geometria algébrica}, I por A. Grothendieck, II por P. Deligne e N. Katz, \emph{Lecture Notes in Mathematics} n\textsuperscript{os}~288, 340, Springer-Verlag, 1972, 1973.}
\endgroup

% PDF page 10: intentionally blank Roman page X.
\clearpage
\thispagestyle{empty}
\null
\clearpage

% Printed pages XI--XII.
\thispagestyle{empty}
\section*{\uline{Índice geral\textsuperscript{(*)}}}

\begingroup
\newcommand{\sgatocentry}[3]{%
  \par\smallskip\noindent
  \begin{minipage}[t]{0.82\linewidth}\textbf{Exposição #1}\end{minipage}%
  \hfill\begin{minipage}[t]{0.12\linewidth}\raggedleft #3\end{minipage}\par
  \noindent\hspace*{2.5em}\begin{minipage}[t]{0.87\linewidth}#2\end{minipage}\par}
\newcommand{\sgatocindex}[2]{%
  \par\smallskip\noindent\hspace*{2.5em}%
  \begin{minipage}[t]{0.72\linewidth}#1\end{minipage}%
  \hfill\begin{minipage}[t]{0.12\linewidth}\raggedleft #2\end{minipage}\par}

\sgatocentry{I}{Complexos dualizantes, por A. Grothendieck, escrito por L. Illusie.}{1}

\sgatocentry{III}{Fórmula de Lefschetz, por A. Grothendieck, escrito por L. Illusie.}{73}

\sgatocentry{III B}{Cálculo dos termos locais, por L. Illusie.}{138}

\sgatocentry{V}{Sistemas projetivos $J$-ádicos, por J.-P. Jouanolou.}{204}

\sgatocentry{VI}{Cohomologia $\ell$-ádica, por J.-P. Jouanolou.}{251}

\sgatocentry{VII}{Cohomologia de alguns esquemas clássicos e teoria cohomológica das classes de Chern, por J.-P. Jouanolou.}{282}

\sgatocentry{VIII}{Grupos de classes das categorias abelianas e trianguladas, complexos perfeitos, por A. Grothendieck, escrito por I. Bucur.}{351}

\sgatocentry{X}{Fórmula de Euler--Poincaré em cohomologia étale, por A. Grothendieck, escrito por I. Bucur.}{372}

\vfill
\noindent\textsuperscript{(*)} Para explicações sobre a ausência das exposições II, IV, IX, XI e XIII, veja a introdução.

\clearpage

\sgatocentry{XII}{Fórmulas de Nielsen--Wecken e de Lefschetz em geometria algébrica, por A. Grothendieck, escrito por I. Bucur.}{407}

\sgatocentry{XIV = XV}{Morfismo de Frobenius e racionalidade das funções $L$, por C. Houzel.}{442}

\sgatocindex{Índice terminológico}{481}
\sgatocindex{Índice de notações}{483}
\endgroup

\bigskip
\[
\begin{tikzcd}[column sep=2.5em,row sep=1.6em]
  & & \mathrm{III} \arrow[r] & \mathrm{III\,B} \arrow[dr] & {} \\
  \mathrm{I,\ VII,\ VIII} \arrow[r] & \mathrm{X} & \mathrm{V} \arrow[r] & \mathrm{VI} \arrow[r] & \mathrm{XIV} \\
  & & \mathrm{III\,B\,5} \arrow[r] & \mathrm{XII} \arrow[ur] & {}
\end{tikzcd}
\]

\clearpage
\pagenumbering{arabic}
\setcounter{page}{1}
\pagestyle{plain}
